% answers_short.tex —— 2009 年·普通高等学校招生考试·辽宁卷·理科 · 简短答案速查
% 由 tools/make_short_answers.py 自 answers.tex 的 \gksolhead/\gkanswer 提取，
% 并与 2009/2009_数学真题_辽宁理_答案要点.md 逐题比对后生成，勿手改（改源 MD 后重跑）。

\begin{gknotice}
\textbf{使用说明}：本节只给最终答案便于快速核对，完整推导见解析版。答案为编写者独立推导，\textbf{非官方标准答案}。
\end{gknotice}

\gkpart{选择题}
\begin{center}\begin{tabular}{ccccccccccccc}
\toprule
题号 & 1 & 2 & 3 & 4 & 5 & 6 & 7 & 8 & 9 & 10 & 11 & 12 \\
\midrule
答案 & B & C & B & B & A & B & D & B & A & C & C & C \\
\bottomrule
\end{tabular}\end{center}

\gkpart{填空题}
\begin{itemize}[leftmargin=2.6em]
  \item[\textbf{13.}] $1013$。
  \item[\textbf{14.}] $\frac13$。
  \item[\textbf{15.}] $4$。
  \item[\textbf{16.}] $9$。
\end{itemize}

\gkpart{解答题}
\begin{itemize}[leftmargin=2.6em]
  \item[\textbf{17.}] $BD=AB$，$BD\approx0.33\ \mathrm{km}$。
  \item[\textbf{18.}] （1）$\frac{\sqrt6}{3}$；（2）证明见解析。
  \item[\textbf{19.}] （1）$P(X=k)=\mathrm C_4^k\left(\frac13\right)^k\left(\frac23\right)^{4-k}$（$k=0,1,2,3,4$）；（2）$\frac7{25}$。
  \item[\textbf{20.}] （1）$\frac{x^2}{4}+\frac{y^2}{3}=1$；（2）定值为 $\frac12$。
  \item[\textbf{21.}] （1）按 $1<a<2$、$a=2$、$a>2$ 分列单调区间；（2）证明见解析。
  \item[\textbf{22.}] 【A】（1）证明见解析；（2）$4\pi$；【B】（1）$x+\sqrt3y=2$，$M(2,0)$，$N\left(\frac{2\sqrt3}{3},\frac\pi2\right)$；（2）$\theta=\frac\pi6$；【C】（1）$\left(-\infty,-\frac32\right]\cup\left[\frac32,+\infty\right)$；（2）$a\leqslant-1$ 或 $a\geqslant3$。
\end{itemize}
